\documentclass{article}

\usepackage{corl_2026} % Use this for the initial submission.
% \usepackage[final]{corl_2026} % Uncomment for the camera-ready ``final'' version.
% \usepackage[preprint]{corl_2026} % Uncomment for pre-prints (e.g., arxiv); This is like ``final'', but will remove the CORL footnote.

% Encoding
\usepackage[utf8]{inputenc}
\usepackage[T1]{fontenc}

% Core math
\usepackage{amsmath, amssymb, amsfonts}
\usepackage{mathtools}
\usepackage{bm}

% Figures and tables
\usepackage{graphicx}
\usepackage{booktabs}
\usepackage{tabularx}
\usepackage{multirow}
\usepackage{makecell}
\usepackage{colortbl}
\usepackage{siunitx}
\usepackage{rotating}
\usepackage{wrapfig}

% Formatting
\usepackage{microtype}
\usepackage[table]{xcolor}
\usepackage{nicefrac}
% \usepackage{refcheck}

% Links (load once, near the end)
% \usepackage[bookmarks=true]{hyperref}

% Algorithms and code
\usepackage[ruled]{algorithm2e}
\usepackage{listings}

% Subfigures
\usepackage[caption=false,font=footnotesize]{subfig}

% Misc
\usepackage{multicol}
\usepackage{dsfont}
\usepackage{tabularray}

% ── Math macros (inline, no external macros.tex required) ─────────────────────
% \newcommand{\reals}{\mathbb{R}}
% \newcommand{\naturals}{\mathbb{N}}
% \newcommand{\expectation}{\mathbb{E}}
% \newcommand{\wspace}{\mathcal{W}}
% \newcommand{\trans}{\mathcal{T}}
% \newcommand{\simulate}{\textsc{Simulate}}
% \newcommand{\config}{\mathcal{C}}
% \newcommand{\states}{\mathcal{S}}
% \newcommand{\rvisit}{R_{\mathrm{visit}}}
% \newcommand{\rterm}{R_{\mathrm{term}}}
% \newcommand{\rpath}{R_{\mathrm{path}}}
% ─────────────────────────────────────────────────────────────────────────────

% Formatting and General
\newcommand{\etal}{\textit{et al. }}
\newcommand{\solution}{\textbf{Solution:}}
\newcommand{\qed}{\hfill \ensuremath{\Box}}
\definecolor{darkgreen}{rgb}{0.0, 0.5, 0.0} 

% Math Sets and Symbols
\newcommand{\reals}{\ensuremath{\mathbb{R}}}
\newcommand{\naturals}{\ensuremath{\mathbb{N}}}
\newcommand{\integers}{\ensuremath{\mathbb{Z}}}
\newcommand{\expectation}{\ensuremath{\mathbb{E}}}
\newcommand{\D}{\mathrm{d}}
\newcommand{\reg}{\ensuremath{\emph{R}}}

% Spaces and Sets
\newcommand{\config}{\mathcal{Q}}
\newcommand{\wspace}{\mathcal{W}}
\newcommand{\obst}{\mathcal{O}}
\newcommand{\goals}{\mathcal{G}}
\newcommand{\states}{\mathcal{S}}
\newcommand{\trans}{\ensuremath{\mathcal{T}}}

% Rewards
\newcommand{\rterm}{R_{\mathrm{term}}}
\newcommand{\rspray}{R_{\mathrm{spr}}}
\newcommand{\rpath}{R_{\mathrm{path}}}
\newcommand{\rvisit}{R_{\mathrm{visit}}}

% Theorem Environments
\newtheorem{problem}{Problem}
\newtheorem{remark}{Remark}
\newtheorem{definition}{Definition}
\newtheorem{proposition}{Proposition}
\newtheorem{lemma}{Lemma}
\newtheorem{theorem}{Theorem}

% Paired Delimiters
\DeclarePairedDelimiter\ceil{\lceil}{\rceil}
\DeclarePairedDelimiter\floor{\lfloor}{\rfloor}

% Operators
\DeclareMathOperator*{\argmax}{arg\,max}
\DeclareMathOperator{\simulate}{SIM}
\DeclareMathOperator{\init}{init}
\DeclareMathOperator{\free}{free}
\DeclareMathOperator{\new}{new}
\DeclareMathOperator{\pos}{POS}
\DeclareMathOperator{\goalf}{GS}
\DeclareMathOperator{\visited}{\mathbb{V}}

% Fallback definitions for general-formulation macros (no-op if already defined in macros.tex)
\providecommand{\obst}{\mathcal{O}}
\providecommand{\free}{\mathrm{free}}
\providecommand{\obs}{\mathrm{obs}}
\providecommand{\col}{\mathrm{col}}

% Fallback definitions for general-formulation macros (no-op if already defined in macros.tex)
\providecommand{\obst}{\mathcal{O}}
\providecommand{\free}{\mathrm{free}}
\providecommand{\obs}{\mathrm{obs}}
\providecommand{\col}{\mathrm{col}}

\title{Learning to Navigate in Stochastic and Uncertain Environments with Different Robot Dynamics}

% The \author macro works with any number of authors. There are two
% commands used to separate the names and addresses of multiple
% authors: \And and \AND.
%
% Using \And between authors leaves it to LaTeX to determine where to
% break the lines. Using \AND forces a line break at that point. So,
% if LaTeX puts 3 of 4 authors names on the first line, and the last
% on the second line, try using \AND instead of \And before the third
% author name.

% NOTE: authors will be visible only in the camera-ready and preprint versions (i.e., when using the option 'final' or 'preprint').
% 	For the initial submission the authors will be anonymized.

\author{
  Jane E.~Doe\\
  Department of Electrical Engineering and Computer Sciences\\
  University of California Berkeley
  United States\\
  \texttt{janedoe@berkeley.edu} \\
  %% examples of more authors
  %% \And
  %% Coauthor \\
  %% Affiliation \\
  %% Address \\
  %% \texttt{email} \\
  %% \AND
  %% Coauthor \\
  %% Affiliation \\
  %% Address \\
  %% \texttt{email} \\
  %% \And
  %% Coauthor \\
  %% Affiliation \\
  %% Address \\
  %% \texttt{email} \\
  %% \And
  %% Coauthor \\
  %% Affiliation \\
  %% Address \\
  %% \texttt{email} \\
}

\begin{document}
\maketitle
\thispagestyle{empty}
\pagestyle{empty}

%===============================================================================

\begin{abstract}
Autonomous multi-robot systems operating under real-world conditions must be robust to wind disturbances, actuator noise, and sensor imperfections, yet existing multi-robot navigation benchmarks typically assume deterministic dynamics and fixed robot platforms, limiting their real-world applicability. We present a formally specified, dynamics-agnostic framework for learning-based multi-robot navigation under stochastic wind, actuation noise, and observation uncertainty. We formulate the navigation task as a constrained stochastic optimization over target coverage and prove its equivalence to a Partially Observable Markov Decision Process (POMDP), yielding a theoretical guarantee absent in prior benchmarks. We instantiate the framework with two physically distinct platforms---UAVs with Newtonian point-mass dynamics and wheeled robots with bicycle kinematics---sharing the same task structure, uncertainty modes, and domain randomization protocol. We benchmark both on-policy and off-policy deep RL algorithms with systematic ablations over reward function, observation space, uncertainty sources, and domain randomization (DR), and quantify the UAV sim-to-real observation gap in CoppeliaSim under position noise, wind-estimate latency, and communication dropout.
\end{abstract}


\section{Introduction}

Autonomous multi-robot systems must reliably navigate complex, unstructured environments where stochastic disturbances---wind gusts, imprecise actuators, and noisy sensors---are unavoidable. Reinforcement learning (RL) has emerged as a compelling approach to multi-robot navigation \citep{chen2017icra, everett2018iros, fan2020ijrr}. However, progress in applying RL to multi-robot systems with continuous control and real-world uncertainties has been limited, and the design of rigorous evaluation environments remains an open challenge \citep{henderson2018deep, reda2020learning, dulacArnold2021}. In prior deep RL works on multi-robot navigation \citep{chen2017icra, everett2018iros, fan2020ijrr, long2017iral, han2022iral}, formal guarantees linking the navigation objective to the optimal policy are absent, and robot dynamics are fixed to a single platform with no stochasticity or physical uncertainty. Without a formal link between the environment's reward signal and the underlying navigation objective, it is difficult to interpret whether learned policies actually optimize the intended task \citep{reda2020learning, dulac2020empirical}.

In this paper, we center the question: \emph{how do standard deep RL methods behave when multi-robot navigation is deliberately made stochastic and uncertain, and when the robot dynamics differ substantially between platforms?} This paper introduces a dynamics-agnostic stochastic navigation framework and instantiate it with two RL environments---one for UAVs and one for wheeled robots---that share the same task structure, algorithmic pipeline, and experimental ablation protocol, yet expose completely different motion dynamics. We formulate the cooperative target-coverage task---where a team of robots must visit all goal locations while minimizing path costs and avoiding collisions---as a constrained stochastic optimization problem using a platform-independent robot model and a unified free-workspace representation that supports workspaces with or without static obstacles. We derive the equivalent POMDP and
prove that the optimal POMDP policy solves the underlying optimization problem. Our specific contributions are:
\begin{enumerate}
\item \textbf{Dynamics-agnostic framework.} We present a general constrained optimization formulation for multi-robot target coverage with collision avoidance, built on a platform-independent robot model~\citep{le2021icra} and a unified free-workspace representation that naturally accommodates workspaces with or without static obstacles.

\item \textbf{Optimization-to-POMDP equivalence.} We formally prove that the POMDP's optimal policy solves the constrained navigation optimization problem, providing theoretical grounding absent in prior works \citep{chen2017icra, everett2018iros, fan2020ijrr}.

\item \textbf{Dual-platform stochastic environments.} We release two RL environments instantiating the general framework under shared stochastic wind, actuation noise, observation noise, and domain randomization, enabling direct comparison across platforms under identical experimental conditions.

\item \textbf{Comprehensive benchmarking.} We evaluate six deep RL algorithms (A2C \citep{mnih2016a2c}, ARS \citep{mania2018ars}, TRPO \citep{schulman2015trust}, PPO
\citep{schulman2017ppo}, TQC \citep{kuznetsov2020tqc}, CrossQ \citep{bhatt2024crossq}) and perform systematic ablations over reward function, observation space, uncertainty sources, and domain randomization.

\item \textbf{Sim-to-real gap analysis.} We deploy trained UAV policies in a 3-D CoppeliaSim~\citep{coppeliaSim2013} scene and quantify performance degradation under position noise, wind-estimate latency, and communication dropout, establishing a baseline for transfer research.
\end{enumerate}

\section{Related Work}
\label{sec:recent}
 
Classical planning methods for multi-robot navigation — including A* \citep{goto2003}, RRT/RRT* \citep{rodriguez2006rrt, karaman2011rrts}, CBS \citep{sharon2015cbs, li2019mccbs}, PDDL-based planners \citep{fox2003pddl21, coles2010popf, benton2012optic}, ScottyActivity \citep{fernandez2018scottyactivity}, and sampling-based approaches with dynamics \citep{le2017icaps, le2019ral, le2021icra, bui2022iros} — rely on deterministic models and produce plans that are susceptible to compounding errors under continuous stochastic disturbances. RL addresses this by optimizing expected cumulative reward under the full distribution of dynamics and noise parameters \citep{reda2020learning, dulacArnold2021, ibarz2021ijrr, tang2025}.
 
Deep RL has been applied to multi-robot navigation with increasing scale and fidelity. \citet{bae2019as} used DQN in grid-based settings; \citet{chen2017icra} demonstrated decentralized collision avoidance in continuous space; \citet{everett2018iros} scaled this to $n$-robot teams via LSTMs; \citet{long2017iral} and \citet{fan2020ijrr} further showed coordination of large robot teams in cluttered environments; and \citet{chen2021iral} addressed decentralized task and path planning jointly. However, these works assume deterministic or simplified dynamics: wind disturbances, actuation noise, and heterogeneous robot platforms are not modeled within a single benchmark. MARL algorithms such as VDN \citep{sunehag2018vdn}, QMIX \citep{rashid2018qmix}, MADDPG \citep{lowe2017maddpg}, and MAPPO \citep{yu2022mappo} provide scalable coordination frameworks, while off-policy methods TQC \citep{kuznetsov2020tqc} and CrossQ \citep{bhatt2024crossq} offer high sample efficiency in continuous control. Domain randomization \citep{tobin2017domain, peng2018sim2real} bridges the sim-to-real gap but is typically applied to a single platform without formal task guarantees.
 
On environment design, benchmarks such as Google Research Football \citep{kurach2020google} and CityFlow \citep{zhang2019cityflow}, and simulators such as MuJoCo \citep{todorov2012mujoco} and Isaac Gym \citep{makoviychuk2021isaac}, have advanced RL research \citep{henderson2018deep, reda2020learning, dulac2020empirical} but are not designed for multi-robot navigation under physical uncertainty and provide no formal link between reward and navigation objectives. Our framework addresses all three gaps: it is dynamics-agnostic, supports stochastic disturbances across two physically distinct platforms, and formally proves that the optimal POMDP policy solves the underlying constrained navigation problem. Building on a general robot model for multi-robot motion planning \citep{le2021icra}, our framework is dynamics-agnostic, explicitly incorporates multiple sources of physical uncertainty, and provides a formal proof linking the POMDP reward to the underlying constrained optimization objective.

\section{Preliminaries}

Let $\reals, \reals_{\geq 0}, \naturals$ be the set of real, non-negative real, and natural numbers, respectively. Let $\{x_i\}_{i=1}^n$ denote an ordered sequence of elements $x_1, \ldots, x_n$. The $L^2$ norm $\|.\|_2$ is simplified as $\|.\|$ for convenience. Let $\mathbf{I}_n$ denote the $n \times n$ identity matrix. Let $\mathbf{0}_n$ and $\mathbf{1}_n$ denote $n$-dimensional all-zeros and all-ones vectors. Let $\mathds{1}_C$ denote a boolean variable where $\mathds{1}_C=1$ if boolean expression $C$ is true and $0$ otherwise. Let $\mathbb{P}(E)$ denote the probability of an event $E$. For a set of samples $S$, let $\Delta(S)$ denote probability distributions over $S$. For a discrete random variable $Y$ taking values in a countable set $\mathcal{Y}$ with probability mass function $p(y) = \mathbb{P}(Y=y)$, the expected value (or expectation) is defined as $\mathbb{E}[Y] = \sum_{y \in \mathcal{Y}} y p(y)$. A sample from distribution $\mathcal{X}$ with mean $\mu \in \reals$ and variance $\sigma^2 \in \reals_{\geq 0}$ is written as $x \sim \mathcal{X}(\mu, \sigma^2)$. Let $\mathcal{N}(\mu, \sigma^2)$ denote the univariate normal (Gaussian) distribution. Let $\mathcal{N}(\bm{\mu}, \mathbf{\Sigma})$ denote the multivariate Gaussian distribution with mean vector $\bm{\mu} \in \reals^n$ and covariance matrix $\mathbf{\Sigma} \in \reals^{n \times n}$, where $\mathbf{\Sigma}$ is symmetric positive semi-definite \citep{bishop2006pattern}. In the context of decision-making, a closed-loop policy $\pi$ is defined as a mapping from a state space $S$ to a probability distribution over an action space $A$, denoted $\pi(a \mid s) \in \Delta(A)$. We use the notation and definitions from \citet{choset2005} to formulate the multi-robot navigation problem in the next section.


\section{Problem Formulation}
\label{sec:prob}
Consider a group of $n \in \naturals$ mobile robots $\mathcal{M} = \{ \mathcal{M}_1, \ldots, \mathcal{M}_n \}$ operating in a bounded $d$-dimensional workspace $\wspace \subseteq \reals^d$ that may contain a set of static obstacles $\obst$. Let $\wspace_{O_j}$ denote the region occupied by obstacle $j$. The obstacle-free (free) workspace is $\wspace^{\free} = \wspace \setminus \cup_j \wspace_{O_j}$. The robots cooperate to visit a set of $m \in \naturals$ target locations (or goals) $\goals = \{ g_1, \ldots, g_m \}$. Let $p_j^g \in \wspace$ denote the position of target $j$ and $r_j > 0$ denote its coverage radius. Following \citet{le2021icra}, each robot $\mathcal{M}_i = \{ \mathcal{P}_i, \config_i, \states_i, s_i^{\init}, \mathcal{U}_i, f_i \}$ is described by its shape $\mathcal{P}_i$, configuration space $\config_i$, state space $\states_i$, initial state $s_i^{\init} \in \states_i$, control space $\mathcal{U}_i$, and dynamics function $f_i$. At configuration $q_i \in \config_i$, robot $i$ occupies region $\mathcal{P}_i(q_i) \subset \wspace$. The obstacle-occupied and collision configuration spaces are:
\begin{align}
    \config_i^{\obs} &= \{q_i \in \config_i \mid \exists  j : \mathcal{P}_i(q_i) \cap \wspace_{O_j} \neq \emptyset\}, \label{eq:obs_config} \\
    \config_i^{\col} &= \{q_i \in \config_i \mid \exists  j \neq i : \mathcal{P}_i(q_i) \cap \mathcal{P}_j(q_j) \neq \emptyset\} \cup \config_i^{\obs}. \label{eq:col_config}
\end{align}
Since the robots are moving, let $\config_i^{\col}(t)$ denote the collision configuration for robot $i$ at time $t$. The free configuration space at time $t$ is $\config_i^{\free}(t) = \config_i \setminus \config_i^{\col}(t)$. The state space $\states_i = \config_i \times \mathcal{V}_i$ augments the configuration space with motion parameters $\mathcal{V}_i$ (e.g., velocity, heading, steering angle). The dynamics of each robot satisfy:
\begin{equation}
    \dot{s}_i = f_i(s_i, u_i), \quad s_i \in \states_i,\ u_i \in \mathcal{U}_i. \label{eq:dynamics}
\end{equation}
The dynamics follow a forward Euler discretization of a continuous-time process \citep{feedback2020}. A new state $s_i^{\mathrm{new}} \in \states_i$ reached after a small time step $dt$ is computed by numerically integrating $f_i$:
\begin{equation}
    s_i^{\mathrm{new}} \gets \simulate(s_i, u_i, f_i, dt). \label{eq:simulate}
\end{equation}
The environment experiences a stochastic wind with per-step magnitude $v_w(t) \sim \mathcal{N}(\bar{v}_w, \epsilon_w^2)$ and direction $\theta_w(t) \sim \mathcal{N}(\bar{\theta}_w, \epsilon_\theta^2)$, yielding wind vector $\mathbf{w}(t) = v_w(t) [\cos\theta_w(t),  \sin\theta_w(t)]^T$ where $\bar{v}_w$, $\bar{\theta}_w$ are nominal values and $\epsilon_w^2$, $\epsilon_\theta^2$ are noise variances. How $\mathbf{w}(t)$ changes the dynamics depends on the platform and is specified in Section~\ref{sec:exp:dynamics}. Each target $g_j \in \goals$ carries a binary visit flag $b_j(t) \in \{0, 1\}$, initialized to $0$ and transitioning irreversibly to $1$ once any robot comes within $r_j$ of $p_j^g$:
\begin{equation}
    b_j(t+dt) = \max \Big( b_j(t),  \max_{i} \mathds{1}_{\|p_i(t) - p_j^g\| \leq r_j} \Big). \label{eq:coverage}
\end{equation}
Let $\mathbf{b}(t) = (b_1(t), \ldots, b_m(t)) \in \{0,1\}^m$ denote the full visit flags at time $t$. Let $\mathcal{J} : \states_i \times \mathcal{U}_i \to \reals_{\geq 0}$ be a per-step cost function:
\begin{equation}
    \mathcal{J}(s_i(t), u_i(t)) = \alpha_1 \|u_i(t)\|^2 + \alpha_2 \|v_i(t)\| + \alpha_3, \label{eq:cost}
\end{equation}
where $v_i(t) \in \mathcal{V}_i$ are the motion parameters, $\alpha_1 \|u_i(t)\|^2$ is the energy cost, $\alpha_2 \|v_i(t)\|$ is the movement cost, $\alpha_3$ is the per-step time cost, and $\alpha_1, \alpha_2, \alpha_3 \in \reals_{\geq 0}$ are fixed weights. 

Because the dynamics are subjected to unbounded stochastic noise, strict constraint satisfaction for all possible noise realizations is impossible. Instead, we formulate this as a stochastic optimal control problem over a joint closed-loop policy $\pi$. Let $T$ be the random stopping time when either all targets are visited (success) or a constraint violation occurs (collision or leaving the free workspace). The objective is to find a policy $\pi$ that maximizes the probability of success, and among such policies, minimizes the expected cumulative path cost:
\begin{flalign}
 & \min_{\pi} \mathbb{E}_{\pi} \Big[ \sum_{t=0}^{T} \sum_{i=1}^{n} \mathcal{J}(s_i(t), u_i(t))  \Big|  s_i(0) = s_i^{\init}, \mathbf{b}(0) = \mathbf{0}_m \Big] \label{eq:opt} \\
\text{s.t.,} \nonumber \\
\text{(i) }  &  \sum_{g_j \in \goals} b_j(T) = m,  \label{eq:const1} \\
\text{(ii) } & \forall t \in [0,T], \forall \mathcal{M}_i \in \mathcal{M}: \mathcal{P}_i(q_i(t)) \subset \wspace^{\free}, \label{eq:const2} \\
\text{(iii) }& \forall t \in [0,T], \forall \mathcal{M}_i \neq \mathcal{M}_j: \mathcal{P}_i(q_i(t)) \cap \mathcal{P}_j(q_j(t)) = \emptyset. \label{eq:const3}
\end{flalign}
Here, constraint (i) ensures maximizing the probability of visiting all target locations within time $T$. Constraints (ii) and (iii) define the safe operating conditions for avoiding collisions prior to termination. Since it is a \textit{stochastic} optimization problem, constraint satisfaction is considered with high probability, and guaranteed for deterministic idealization (without noise),. 

\section{Environment Design}
\label{sec:mdp}

The optimization problem can be formulated as a POMDP $M = (S, A, \Omega, O, \trans, R, \gamma)$. The \textbf{State Space} $S = \prod_{i=1}^n \states_i \times \{0,1\}^m$ captures the joint state $s_t = \big(\{s_i(t)\}_{i=1}^n,  \mathbf{b}(t)\big)$, where $\states_i$ is the platform-specific state space (details in Section~\ref{sec:exp:dynamics}). The \textbf{Action Space} $A = \mathcal{U}_1 \times \cdots \times \mathcal{U}_n$, where $\mathcal{U}_i$ is the platform-specific control space. The \textbf{Observation Space} $\Omega$ provides a noisy, partial view of the true state. The visit history $\mathbf{b}(t)$ is compressed into a single scalar via the binary-to-decimal encoding $\phi(\mathbf{b}) = \sum_{j=1}^m b_j 2^{j-1}$. Gaussian noise $\mathcal{N}(\mathbf{0}_c, \sigma_\mathrm{obs}^2 \mathbf{I}_c)$ is added to all observations where $\sigma_\mathrm{obs}^2 \in \reals_{\geq 0}$ is the joint observation variance and $c \in \naturals$ is the observation dimension. Platform-specific observation dimensions for each mode are given in Table~\ref{tab:obs_modes} in Section~\ref{sec:exp:dynamics}. The \textbf{Observation Function} $O : S \to \Delta(\Omega)$ maps each true state to a distribution over noisy observations as described above. The \textbf{Transition Function} $\trans : S \times A \to \Delta(S)$ governs the stochastic dynamics:
\begin{equation}
    \trans(s_t, a_t) = \Big(\simulate(s_1, u_1, f_1, dt), \ldots, \simulate(s_n, u_n, f_n, dt),  b_{t+dt}\Big), \label{eq:transition}
\end{equation}
where $f_i$ incorporates wind and actuation noise, and the visit flags update via Eq.~\eqref{eq:coverage}. The \textbf{Reward Function} $R : S \times A \to \reals$ balances target coverage, collision avoidance, and path efficiency. Let $r_s, r_m, r_L \in \reals_{\geq 0}$ be some small, medium, and very large rewards ($r_s < r_m \ll r_L$). The total reward decomposes as:
\begin{alignat}{2}
    R(s_t, a_t) &= \rvisit(t) + \rterm(t) + \rpath(t), \label{eq:reward}\\
    \rvisit(t)  &= r_m \cdot |\{g_j \in \goals : b_j(t + dt) = 1,  b_j(t) = 0\}|, \label{eq:r_visit} \\
    \rterm(t)   &=
    \begin{cases}
        -r_L, & \text{if } \exists \mathcal{M}_i: q_i \in \config_i^{\col}(t) \\
        +r_L, & \text{else if } \sum_{g_j \in \goals} b_j(t) = m \\
        0,    & \text{otherwise}
    \end{cases} \label{eq:r_term} \\
    \rpath(t)   &= - \sum_{i=1}^n \mathcal{J}(s_i(t), u_i(t))
                   - r_m \sum_{i=1}^n \mathds{1}_{p_i(t) \notin \wspace^{\free}}, \label{eq:r_path}
\end{alignat}
where $\rvisit$ rewards newly visited targets, $\rterm$ is the terminal reward for visiting all locations or exiting on collision, and $\rpath$ penalizes energy, speed, path length, and out-of-free-workspace steps. Finally, the \textbf{Discount Factor} $\gamma \in [0,1)$ incentivizes rapid target coverage. We define the following theorem to describe the relation between the optimization problem and the developed POMDP:

\begin{theorem}
\label{thm:optimization}
In the limit $\gamma \to 1$, the optimal policy $\pi^*$ for the underlying fully-observable MDP of $M$ solves the stochastic optimization problem in Eq.~\eqref{eq:opt}--\eqref{eq:const3} starting at state $s = (\{s_i^{\mathrm{init}}\}_{i=1}^n,  \mathbf{0}_m)$. Conversely, any optimal policy for the optimization problem is an optimal policy for the MDP.
\end{theorem}

\textbf{Proof.} The proof analyzes the underlying MDP where the true state $s_t$ is known, which serves as the theoretical upper bound for the POMDP and establishes the correctness of the reward shaping. For a finite horizon $[0,T]$, the definition of optimal policy $\pi^*$ and optimal action-value function $Q^{\pi^*}$ \citep{sutton2018reinforcement} starting at state $s_0 = (\{s_i^{\mathrm{init}}\}_{i=1}^n,  \mathbf{0}_m)$ are as follows:
\begin{align}
    \pi^*(s_0) &= \argmax_a Q^{\pi^*}(s_0, a), \label{eq:opt_pi} \\
    Q^{\pi^*}(s_0, a_0) &= \max_\pi \expectation_\pi \Big[\sum_{\tau=0}^T \gamma^\tau R(s_\tau, a_\tau)  \mid s_0 = (\{s_i^{\mathrm{init}}\}_{i=1}^n,  \mathbf{0}_m) , a_0\Big] \label{eq:opt_policy} \\
    = \max_\pi &\Big( \expectation_\pi \big[ \sum_{\tau=0}^T\gamma^\tau \rvisit(\tau) \mid . \big] + \expectation_\pi \big[ \sum_{\tau=0}^T\gamma^\tau \rterm(\tau) \mid . \big] + \expectation_\pi \big[ \sum_{\tau=0}^T\gamma^\tau \rpath(\tau) \mid . \big] \Big) \tag*{[From Eq. \eqref{eq:reward} and properties of $\mathbb{E}$]}
\end{align}
where $T$ is the stopping time when the episode terminates (either all targets are visited or a collision occurs). We prove both directions of the equivalence.

\textbf{(Direction $\Rightarrow$) $\pi^*$ solves the optimization problem:}
We proceed by analyzing the components of $Q^\pi(s_0, a_0)$ as $\gamma \to 1$ in two lemmas.

\begin{lemma}
\label{lem:const}
Assuming $r_L \gg \sum_{\tau=0}^T \gamma^\tau (\rvisit(\tau) + \rpath(\tau))$, any optimal policy $\pi^*$ must maximize the probability of satisfying constraints~\textnormal{(i)}, \textnormal{(ii)}, and~\textnormal{(iii)} of Eq.~\eqref{eq:opt}.
\end{lemma}

\textit{Proof.} Because actuation and position noises are unbounded, no policy can eliminate the possibility of a constraint violation with probability exactly~1. We show that reward shaping makes constraint violations overwhelmingly suboptimal.

\paragraph{Constraint (iii) --- no collisions.} Suppose $\pi^*$ induces a collision at some time $t^*$; then $\rterm(t^*) = -r_L$. Here, $r_L \gg \sum_{\tau=0}^T \gamma^\tau (\rvisit(\tau) + \rpath(\tau))$ makes $Q^\pi \approx -r_L$, which is the lower bound. No such policy can be optimal.

\paragraph{Constraint (ii) --- obstacle and workspace containment.} Any policy placing robot $i$ outside $\wspace^{\free}$ (i.e., outside the workspace boundary or into an obstacle region) incurs a per-step penalty $-r_m$ in $\rpath$. Since $\wspace^{\free} = \wspace \setminus \bigcup_j \wspace_{O_j}$, this uniformly penalizes both boundary violations and obstacle penetrations. Keeping $\mathcal{P}_i(q_i) \subset \wspace^{\free}$ sets this term to zero, maximizing $\rpath$ since $r_m > 0$.

\paragraph{Constraint (i) --- all targets visited.} Visiting all $m$ targets yields $\rterm = +r_L$. Since $r_L$ dominates all other accumulated rewards, maximizing $\rterm$ ensures the highest probability for $\sum_{g_j \in \goals} b_j(T) = m$ at termination, satisfying constraint~(i). \hfill$\square$

\begin{lemma}
\label{lem:cost}
Among all policies satisfying constraints~\textnormal{(i)--(iii)}, maximizing $Q^\pi$ is equivalent to minimizing the cost objective of Eq.~\eqref{eq:opt}.
\end{lemma}

\textit{Proof.} By Lemma~\ref{lem:const} and constraint~(i), $\rterm(\tau) = +r_L$ at the terminal step $T$ and $0$ otherwise, so $\expectation_\pi[\sum_{\tau=0}^T\gamma^\tau \rterm(\tau) \mid .] = \expectation_\pi[\gamma^T r_L \mid .] \approx r_L$ for $\gamma \to 1$.

For the visit reward, each target $j$ is visited exactly once by definition, contributing $r_m$ at the step it is first reached. Therefore:
\begin{equation}
    \expectation_\pi \Big[\sum_{\tau=0}^T \gamma^\tau \rvisit(\tau) \mid .\Big]
    = r_m \sum_{j=1}^m \expectation[\gamma^{t_j} \mid .]
    \approx r_m m, \label{eq:visit_const}
\end{equation}
where $t_j$ is the step when target $j$ is visited, and the approximation holds for $\gamma \to 1$. This quantity is approximately constant (policy-independent for constraint-satisfying policies). By constraint~(ii), $\mathcal{P}_i(q_i(\tau)) \subset \wspace^{\free}$ for all $i, \tau$, so $\mathds{1}_{p_i \notin \wspace^{\free}} = 0$. Hence:
\begin{equation}
    \expectation_\pi \Big[\sum_{\tau=0}^T\gamma^\tau \rpath(\tau) \mid . \Big]
    \approx -\expectation_\pi \Big[\sum_{\tau=0}^T\gamma^\tau \sum_{i=1}^n \mathcal{J}(s_i(\tau), u_i(\tau)) \mid . \Big], \label{eq:path_const}
\end{equation}
From Eq.~\eqref{eq:opt_policy}:
\begin{flalign}
Q^{\pi^*}(s_0, a_0) &\approx \max_\pi \Big( r_L + r_m m - \expectation_\pi \Big[\sum_{\tau=0}^T\gamma^\tau \sum_i \mathcal{J}(s_i(\tau), u_i(\tau))\Big] \Big) \nonumber \\
&= r_L + r_m m - \min_\pi \expectation_\pi \Big[\sum_{\tau=0}^T\gamma^\tau \sum_i \mathcal{J}(s_i(\tau), u_i(\tau))\Big]. \tag*{[$\max_a(C-f(a)) = C -\min_a f(a)$]}
\end{flalign}
Here,  $r_L + r_m m$ is approximately constant (policy-independent) for policies with identical probabilities. Consequently, maximizing $Q^{\pi^*}$ is precisely equivalent to minimizing the expected discounted cost in Eq.~\eqref{eq:opt}, proving the lemma. Therefore, by Lemmas~\ref{lem:const} and~\ref{lem:cost}, $\pi^*$ solves Eq.~\eqref{eq:opt}.

\textbf{(Direction $\Leftarrow$) The optimization solution induces $\pi^*$:} Let $\pi_{\mathrm{opt}}$ be an optimal closed-loop policy for the stochastic optimization problem in Eqs.~\eqref{eq:opt}--\eqref{eq:const3}. By definition, $\pi_{\mathrm{opt}}$ maximizes the probability of satisfying all three constraints, and among such policies, minimizes the expected path cost. Suppose for contradiction that there exists another policy $\pi'$ with $Q^{\pi'} > Q^{\pi_{\mathrm{opt}}}$. Because $r_L$ dominates the reward function (and the action-value function $Q$), $\pi'$ cannot have a lower success probability than $\pi_{\mathrm{opt}}$. If $\pi'$ has the same success probability, then by Lemma~\ref{lem:cost}, $Q^{\pi'} > Q^{\pi_{\mathrm{opt}}}$ implies:
\[
\mathbb{E}_{\pi'} \Big[\sum_{t=0}^T \sum_{i=1}^n \mathcal{J} \big(s_i(t), u_i(t)\big) \mid . \Big]
< \mathbb{E}_{\pi_{\mathrm{opt}}} \Big[\sum_{t=0}^T \sum_{i=1}^n \mathcal{J} \big(s_i(t), u_i(t)\big) \mid . \Big],
\]
which contradicts the optimality of $\pi_{\mathrm{opt}}$ for Eq.~\eqref{eq:opt}. Therefore, $\pi_{\mathrm{opt}}$ must also be an optimal policy for the MDP. Hence, Theorem~\ref{thm:optimization} holds in both directions. \hfill$\square$

\section{Experiments}
\label{sec:exp}

We present a comprehensive experimental evaluation to measure the following: performance of RL methods on both platforms, sensitivity to design choices, out-of-distribution (OOD) robustness under domain shift, and the simulation-to-real observation gap for the UAV platform. Experimental evaluations are conducted for two groups of robots with different dynamics, instantiating the general framework of Section~\ref{sec:prob} as follows:

\paragraph{UAV dynamics (holonomic).}
Each UAV $\mathcal{M}_i$ has state $s_i = (p_i, v_i) \in \reals^4$ where $p_i \in \wspace$ is its 2-D position and $v_i \in [v_{\min}, v_{\max}]^2$ is its velocity vector. The control input $u_i \in [-1,1]^2$ specifies normalized 2-D thrust to change its velocity. Wind acts as an aerodynamic force that modifies velocity each step:
\begin{align}
    v_i(t+dt) &= \min \Big( \max \Big( v_i(t) + \tfrac{\kappa \mu_i(t) + \epsilon_a}{W_i} dt,  v_{\min} \Big),  v_{\max}\Big), \label{eq:uav_vel} \\
    p_i(t+dt) &= p_i(t) + \big(v_i(t) + \mathbf{w}(t)\big) dt, \label{eq:uav_pos}
\end{align}
where $\kappa > 0$ is the thrust power coefficient, $W_i > 0$ is the UAV mass, $\epsilon_a \sim \mathcal{N}(\mathbf{0}_2, \sigma_a^2 \mathbf{I}_2)$ is the actuation noise, and $\min, \max$ functions are applied element-wise. The UAV workspace $\wspace$ is a convex polygon without obstacles ($\obst = \emptyset$, $\wspace^{\free} = \wspace$).

\paragraph{Wheeled dynamics (non-holonomic).}
Each wheeled robot $\mathcal{M}_i$ has state $s_i = (x_i, y_i, \theta_i, v_i, \delta_i) \in \reals^5$, where $(x_i, y_i)$ is its 2-D position, $\theta_i$ is its heading angle, $v_i \in [v_{\min}, v_{\max}]$ is its longitudinal speed, and $\delta_i \in [\delta_{\min}, \delta_{\max}]$ is the front steering angle. The control input $u_i = (a_i, \dot{\delta}_i) \in [-1,1]^2$ specifies a normalized acceleration $a_i$ and steering rate $\dot{\delta}_i$. Unlike the UAV, wind acts on the wheeled robot as a lateral position drift rather than a velocity perturbation---reflecting the fact that ground vehicles are less susceptible to aerodynamic forces but can still be displaced by strong crosswinds:
\begin{align}
    v_i(t+dt) &= \min \big( \max \big( v_i(t) + a_s a_i dt,  v_{\min} \big),  v_{\max}\big), \label{eq:whl_vel} \\
    \delta_i(t+dt) &= \min \big( \max \big( \delta_i(t) + \dot{\delta}_i,  \delta_{\min} \big),  \delta_{\max}\big), \label{eq:whl_steer} \\
    \theta_i(t+dt) &= \theta_i(t) + \tfrac{v_i(t) \tan(\delta_i(t))}{\ell_i} dt, \label{eq:whl_head} \\
    \begin{bmatrix} x_i(t+dt) \\ y_i(t+dt) \end{bmatrix} &=
    \begin{bmatrix} x_i(t) \\ y_i(t) \end{bmatrix}
    + \Big( v_i(t) \begin{bmatrix}\cos\theta_i(t)\\\sin\theta_i(t)\end{bmatrix} + \mathbf{w}(t) \Big) dt, \label{eq:whl_pos}
\end{align}
where $a_s > 0$ is an acceleration scaling factor and $\ell_i > 0$ is the robot's wheelbase. The wheeled workspace is a rectangle $[0, W]\times[0, H]$ containing static polygon obstacles $\obst$. Table~\ref{tab:obs_modes} summarizes the observation space composition for each platform and mode.

\begin{table}[tbp]
\centering
\footnotesize
\caption{\small Observation space composition for each platform and mode ($N$~=~number of robots). UAV kinematics: position~$(2N)$ + velocity~$(2N)$. Wheeled kinematics: $(x,y,\theta,v,\delta)$~$(5N)$.}
\label{tab:obs_modes}
\begin{tabular}{l|ll}
    \toprule
    \textbf{Mode} & \textbf{UAV ($\dim$)} & \textbf{Wheeled ($\dim$)} \\
    \midrule
    \texttt{full}        & position + velocity + $\phi(\mathbf{b})$ ($4N+1$) & full kinematics\ + $\phi(\mathbf{b})$ ($5N+1$) \\
    \texttt{no\_pos}     & $\phi(\mathbf{b})$ only ($1$) & $(\theta,v,\delta)$ per robot + $\phi(\mathbf{b})$ ($3N+1$) \\
    \texttt{no\_vis\_hist} & position + velocity ($4N$) & full kinematics\ only ($5N$) \\
    \texttt{pos\_only}   & position only ($2N$) & position only ($2N$) \\
    \bottomrule
\end{tabular}
\vspace{-6pt}
\end{table}

% ─────────────────────────────────────────────────────────────────────────────
\subsection{Setup}
\label{sec:exp:setup}
% ─────────────────────────────────────────────────────────────────────────────

We compare four on-policy (A2C \citep{mnih2016a2c}, ARS \citep{mania2018ars}, PPO \citep{schulman2017ppo}, TRPO \citep{schulman2015trust}) and two off-policy (CrossQ \citep{bhatt2024crossq}, TQC \citep{kuznetsov2020tqc}) algorithms via Stable-Baselines3 \citep{stable-baselines3}. We performed the experiments in two settings: one using the default hyperparameters given in Stable-Baselines3 \citep{stable-baselines3}, and another using a set of tuned hyperparameters selected by Optuna \cite{optuna_2019}. For both robot platforms, we vary the number of robots $N \in \{2, 3, 4, 5\}$, starting positions, target locations, and use 10 environmental variations. For the \textbf{UAV} platform, we define the workspace as convex polygons. For the \textbf{Wheeled} platform, we use rectangular workspaces with distinct obstacle arrangements. Every algorithm--configuration pair is trained for $2 \times 10^6$ timesteps with 3 seeds $(42, 123, 9999)$. The interquartile mean (IQM) \citep{agarwal2021deep} is reported as the primary metric. All experiments were conducted on an HPC cluster with GPUs: NVIDIA\textsuperscript{\textregistered} GeForce RTX 1080Ti, 2080Ti, 3090, L4, L40s, A4000, A4500, A6000; CPUs: Intel\textsuperscript{\textregistered} Xeon\textsuperscript{\textregistered} E5-2680v3, E5-2683v4, Gold 6130 (16--48 cores); and 4--16 GB RAM per node. Simulations were performed on a local desktop with AMD Ryzen\texttrademark\ Threadripper 1950X (16-core) CPU, 64 GB RAM, and NVIDIA\textsuperscript{\textregistered} GeForce RTX 2080 GPU. Each run took approximately 1--6 hours.

% ─── Table 2: Main Results (UAV + Wheeled, unified) ──────────────────────────
\begin{table*}[htbp]
\centering
\footnotesize
\caption{IQM of episodic reward (3 seeds $\times$ 10 configurations $\times$ 50 episodes) under default and tuned hyperparameters. \textbf{(a) UAV platform} (top); \textbf{(b) Wheeled platform} (bottom). $\dagger$~=~best per column. ($-$~=~placeholder; to be filled after experiments complete.)}
\label{tab:main_results}
\resizebox{0.8\linewidth}{!}{%
\begin{tabular}{l|cccc|cccc}
    \toprule
    \multirow{2}{*}{\textbf{Alg.}}
        & \multicolumn{4}{c|}{\textbf{Default HP}}
        & \multicolumn{4}{c}{\textbf{Tuned HP}} \\
    \cmidrule(lr){2-5}\cmidrule(lr){6-9}
        & $N{=}2$ & $N{=}3$ & $N{=}4$ & $N{=}5$
        & $N{=}2$ & $N{=}3$ & $N{=}4$ & $N{=}5$ \\
    \midrule
    \multicolumn{9}{c}{\textbf{(a) UAV Platform}} \\
    \midrule
    A2C    & $-2432.8$  & $-646.6$   & $-1203.3$  & $-1449.6$  & $-2925.3$                   & $-442.0$                    & $510.9$                     & $650.8$                     \\
    ARS    & $-8777.7$  & $-5232.2$  & $-4777.9$  & $-4015.4$  & $-6062.9$                   & $-4693.0$                   & $-4040.1$                   & $-3512.5$                   \\
    PPO    & $4959.8$   & $4422.2$   & $1804.8$   & $1073.3$   & $-2516.4$                   & $-1033.0$                   & $-962.0$                    & $-469.2$                    \\
    TQC    & $3187.5$   & $4124.1$   & $3649.1$   & $667.8$    & $\mathbf{10751.7}^\dagger$  & $\mathbf{10718.5}^\dagger$  & $\mathbf{10716.6}^\dagger$  & $\mathbf{10748.6}^\dagger$  \\
    TRPO   & $9464.0$   & $8839.8$   & $6711.2$   & $6536.0$   & $8949.7$                    & $7252.3$                    & $5634.0$                    & $4590.5$                    \\
    CrossQ & $\mathbf{10995.6}^\dagger$ & $\mathbf{11097.9}^\dagger$ & $\mathbf{11157.6}^\dagger$ & $\mathbf{11117.1}^\dagger$ & $1827.2$ & $1895.8$ & $-460.1$ & $-623.4$ \\
    \midrule
    \multicolumn{9}{c}{\textbf{(b) Wheeled Platform}} \\
    \midrule
    A2C    & $-27779.9$  & $-81538.2$  & $-100413.8$ & $-100540.6$ & $-13940.8$                  & $-71489.2$                  & $-98592.8$                  & $-99903.1$                  \\
    ARS    & $-23723.5$  & $-64314.0$  & $-89875.1$  & $-95695.0$  & $-25393.6$                  & $-57089.1$                  & $-73826.2$                  & $-88579.8$                  \\
    PPO    & $-12376.7$  & $-24074.6$  & $-35626.6$  & $-70088.4$  & $-13903.9$                  & $-29746.2$                  & $-60104.0$                  & $-91052.2$                  \\
    TQC    & $-12307.8$  & $-12473.4$  & $-8968.2$   & $-6552.6$   & $-12314.7$                  & $\mathbf{-12289.8}^\dagger$ & $\mathbf{-10106.2}^\dagger$ & $\mathbf{-8534.5}^\dagger$  \\
    TRPO   & $-12281.6$  & $-17719.0$  & $-35364.8$  & $-69802.6$  & $-12139.0$                  & $-14564.0$                  & $-27543.2$                  & $-73135.9$                  \\
    CrossQ & $-7020.3$   & $-11122.4$  & $-4705.0$   & $-5292.2$  & $1827.2$ & $1895.8$ & $-460.1$ & $-623.4$  \\
    \bottomrule
\end{tabular}}
\vspace{-6pt}
\end{table*}

% ─────────────────────────────────────────────────────────────────────────────
\subsection{Main Results}
\label{sec:exp:main}
% ─────────────────────────────────────────────────────────────────────────────

Table~\ref{tab:main_results} compares all six algorithms on both platforms under default and Optuna-tuned hyperparameters, evaluated with 50 episodes per trained model. \textit{(Results and analysis to be added once experiments are complete.)} For reference, an episode is considered successful when all $m$ targets are visited before the step limit, which activates the $+r_L$ terminal bonus in Eq.~\eqref{eq:r_term}.

% ─── Figures (learning curves) ───────────────────────────────────────────────
% \begin{figure}[tbp]
%     \centering
%     \includegraphics[width=\linewidth]{figures/combined_learning_curves.png}
%     \caption{Learning curves --- Default (top row) vs.\ tuned (bottom row) hyperparameters. Left panels: UAV platform. Right panels: Wheeled platform. Dotted lines indicate the reward threshold for a successful policy.}
%     \label{fig:curves}
%     \vspace{-12pt}
% \end{figure}

% ─────────────────────────────────────────────────────────────────────────────
\subsection{Ablation Study}
\label{sec:exp:ablation}
% ─────────────────────────────────────────────────────────────────────────────

We ablate reward terms, observation spaces, and uncertainty parameters using CrossQ ($N=3$, 3 seeds, 10 variations, 50 episodes) on both platforms. For \textbf{Reward ablations}, we compare the IQM of episode lengths, where episodes terminate either on complete target coverage (all $m$ targets visited) or on a collision. Removing $\rterm$ prevents the policy from learning collision avoidance and coverage incentives. Removing $\rpath$ eliminates energy and path-efficiency penalties; because the discount factor $\gamma$ still implicitly penalizes slow episodes, the impact is expected to be moderate. For \textbf{Observation ablations}, we evaluate the effect of withholding kinematic information (\texttt{no\_pos}), visit history (\texttt{no\_vis\_hist}), or both (retaining \texttt{pos\_only}). We hypothesize that removing position information is catastrophic since robots cannot navigate, and that withholding visit history has a milder effect as robots can re-discover targets. For \textbf{Uncertainty ablations}, we isolate the contributions of wind noise (\texttt{wind\_only}), actuation noise (\texttt{act\_only}), neither (\texttt{deterministic}), or both (\texttt{full}). This identifies which noise source is the primary driver of performance degradation. Table~\ref{tab:ablations_and_dr} summarizes the results.

% ─── Table 3: Ablations + DR ─────────────────────────────────────────────────
\begin{table*}[tbp]
\centering
\small
\caption{\small Ablation and domain randomization results for $N=3$ using CrossQ, over 3 seeds $\times$ 10 configurations $\times$ 50 episodes. Left three blocks: ablation experiments. Right block: DR experiments. $\dagger$~=~best per block/column. UAV~(U) and Wheeled~(W) results shown per row.}
\label{tab:ablations_and_dr}
\resizebox{\textwidth}{!}{%
\begin{tabular}{lcc|lcc|lcc|lccc}
    \toprule
    \multicolumn{3}{c|}{\textbf{Reward ablation}}
    & \multicolumn{3}{c|}{\textbf{Observation ablation}}
    & \multicolumn{3}{c|}{\textbf{Uncertainty}}
    & \multicolumn{4}{c}{\textbf{Domain Randomization}} \\
    \midrule
    Condition & \makecell{Ep.\ len \\ IQM (U)} & \makecell{Ep.\ len \\ IQM (W)}
    & Condition & \makecell{Reward \\ IQM (U)} & \makecell{Reward \\ IQM (W)}
    & Condition & \makecell{Reward \\ IQM (U)} & \makecell{Reward \\ IQM (W)}
    & Condition & \makecell{In-dist.\ \\ IQM (U)} & \makecell{OOD \\ IQM (U)} & CVaR$_{0.1}$ \\
    \midrule
    Full          & $-$ & $-$ & Full          & $-$ & $-$ & Full          & $-$ & $-$ & No DR   & $-$ & $-$ & $-$ \\
    No $\rterm$   & $-$ & $-$ & No pos.       & $-$ & $-$ & Wind only     & $-$ & $-$ & Wind DR & $-$ & $-$ & $-$ \\
    No $\rpath$   & $-$ & $-$ & No visit hist.& $-$ & $-$ & Act.\ only    & $-$ & $-$ & Full DR & $-$ & $-$ & $-$ \\
                  &     &     & Pos.\ only    & $-$ & $-$ & Deterministic & $-$ & $-$ &         &     &     &     \\
    \bottomrule
\end{tabular}}
\vspace{-12pt}
\end{table*}

% ─────────────────────────────────────────────────────────────────────────────
\subsection{Domain Randomization}
\label{sec:exp:dr}
% ─────────────────────────────────────────────────────────────────────────────

We implement domain randomization (DR) \citep{tobin2017domain, peng2018sim2real} by re-sampling physical parameters per training episode under three conditions: (A)~No DR, (B)~Wind DR (randomize wind speed and direction uniformly over $[0, 1] \text{m/s}$), and (C)~Full DR (additionally randomize actuation noise std, robot mass/thrust for UAV, or max speed/steer for wheeled). We evaluate on in-distribution ($\text{wind} \in [0, 0.5] \text{m/s}$) and OOD ($\text{wind} \in (0.5, 2.0] \text{m/s}$) test conditions (right block of Table~\ref{tab:ablations_and_dr}). We additionally evaluate the conditional value-at-risk CVaR$_{0.1}$ \citep{tamar2015cvar} to measure worst-case robustness. We hypothesize that wind DR substantially improves OOD performance while full DR further reduces tail risk at some cost to mean performance.
\vspace{-4pt}

% ─────────────────────────────────────────────────────────────────────────────
\subsection{Simulation and Sim-to-Real Analysis}
\label{sec:exp:s2r}
% ─────────────────────────────────────────────────────────────────────────────

\begin{wraptable}{r}{7cm}
\vspace{-1.8em}
\centering
\caption{\small Sim-to-real observation-gap analysis (CrossQ, UAV, $N = 3$, configuration~1, 50 eval eps). $\Delta$IQM: relative change vs.\ no-perturbation baseline.}
\label{tab:obs_gap}
\setlength{\tabcolsep}{4pt}
\small
\begin{tabular}{lcc}
    \toprule
    \textbf{Perturbation} & \textbf{IQM} & $\Delta$\textbf{IQM (\%)} \\
    \midrule
    None (baseline)                      & $-$ & $0.0$   \\
    GPS noise ($\sigma = 0.05$ m)     & $-$ & $-$     \\
    Wind latency (3 steps)               & $-$ & $-$     \\
    Comm.\ loss (10\%)                   & $-$ & $-$     \\
    All combined                         & $-$ & $-$     \\
    \bottomrule
\end{tabular}
\vspace{-0.8em}
\end{wraptable}

We deploy trained UAV policies in a 3-D CoppeliaSim~\citep{coppeliaSim2013} agricultural scene with PyBullet~\citep{coumans2021pybullet} physics connected via ZeroMQ~\citep{ZeroMQ}. Table~\ref{tab:obs_gap} quantifies IQM degradation under inference-time sensor imperfections applied to the CrossQ policy trained with default hyperparameters ($N = 3$, configuration~1, full DR). Four perturbation conditions are evaluated: GPS position noise ($\sigma = 0.05$ m), wind estimate latency (3 steps), communication packet loss (10\%), and all combined. Wind estimate latency is expected to be the dominant transfer risk, as the UAV policy relies on timely wind information to compensate trajectory deviations. GPS position noise should have a smaller individual effect since the coverage reward is spatially smooth. Applying all perturbations simultaneously is expected to cause the largest collapse in performance, consistent with transfer gaps reported for UAV sim-to-real tasks \citep{peng2018sim2real, akkaya2019solving}. Fig.~\ref{fig:sim} illustrates navigation paths for two UAV environment configurations in the CoppeliaSim scene.

% \begin{figure}[tbp]
%     \centering
%     \subfloat[Config.~1 initial]{\includegraphics[width=0.247\textwidth, height=0.17\textwidth]{figures/simulations/env1_init.png}\label{fig:env1_init}}
%     \hfil
%     \subfloat[Config.~1 final]{\includegraphics[width=0.247\textwidth, height=0.17\textwidth]{figures/simulations/env1_v1_best.png}\label{fig:env1_best}}
%     \hfil
%     \subfloat[Config.~3 initial]{\includegraphics[width=0.247\textwidth, height=0.17\textwidth]{figures/simulations/env3_v1_init.png}\label{fig:env3_init}}
%     \hfil
%     \subfloat[Config.~3 final]{\includegraphics[width=0.247\textwidth, height=0.17\textwidth]{figures/simulations/env3_v1_best.png}\label{fig:env3_best}}
%     \caption{\small CoppeliaSim simulation of UAV navigation for configurations~1 (a, b) and~3 (c, d) with $N=3$ robots. Cyan markers indicate unvisited target locations; dark grey markers indicate visited targets; colored lines trace robot navigation paths.}
%     \label{fig:sim}
%     \vspace{-12pt}
% \end{figure}

\section{Conclusion}
\label{sec:con}

We introduced a dynamics-agnostic framework and dual-platform stochastic benchmark for multi-robot navigation, formally linking the POMDP reward to the optimization objective. Future work includes heterogeneous fleets, decentralized observations, and real-world hardware deployment.


%===============================================================================
\clearpage
% The acknowledgments are automatically included only in the final and preprint versions of the paper.
\acknowledgments{If a paper is accepted, the final camera-ready version will (and probably should) include acknowledgments. All acknowledgments go at the end of the paper, including thanks to reviewers who gave useful comments, to colleagues who contributed to the ideas, and to funding agencies and corporate sponsors that provided financial support.}
%===============================================================================

% no \bibliographystyle is required, since the corl style is automatically used.
\bibliography{ref}  % .bib

\end{document}
